\documentclass[11pt]{article}
\usepackage[margin=1in]{geometry}
\usepackage[T1]{fontenc}
\usepackage[utf8]{inputenc}
\usepackage{lmodern}
\usepackage{amsmath,amssymb}
\usepackage{graphicx}
\usepackage{hyperref}
\usepackage{enumitem}
\usepackage{booktabs}
\setlength{\parskip}{0.5em}
\setlength{\parindent}{0pt}

% Figures are referenced relative to fbc/neural_network/ (compile from that directory).
\graphicspath{{bayes/results/}}

\title{Bayesian Last-Layer Uncertainty in a Neural Feedback Controller:\\
Comparing Random-Walk and Gradient-Aware MCMC}
\author{Arash Mehrabi}
\date{}

\begin{document}
\maketitle

\begin{abstract}
Learned neural controllers emit a single action per state and give no signal when they operate
outside their training distribution. We make the final linear layer of a trained 7-DOF feedback
controller (DNFC) Bayesian and ask how the choice of MCMC sampler shapes the resulting uncertainty.
On the resulting $\approx\!2.7\times10^{3}$-dimensional last-layer posterior we compare a gradient-free
random-walk Metropolis--Hastings (RWMH) sampler against gradient-aware stochastic gradient Langevin
dynamics (SGLD), using split-$\hat R$/ESS diagnostics, a closed-form Gaussian reference posterior, and
results averaged over five independently trained controllers. RWMH tunes to the optimal $23.4\%$
acceptance yet fails to mix ($\mathrm{ESS}\approx5$ from $8000$ draws), exactly as its $1/d$ dimension
scaling predicts; SGLD mixes essentially dimension-independently and is $33\times\!\pm\!15$ more
efficient per unit of compute. We find that posterior \emph{conditioning}---set by the observation-noise
scale $\sigma_a$---governs feasibility for both samplers, that tempering does not help (it rescales
posterior width but not conditioning), and that the predictive uncertainty rises monotonically under
genuine target perturbation, with SGLD recovering the exact posterior's spread far more faithfully than
RWMH. The dataset's nominal ``extrapolation'' split, by contrast, is \emph{not} actually
out-of-distribution (detection AUROC $\approx 0.5$ even for the exact posterior), so we demonstrate the
shift signal through controlled perturbation instead. The practical takeaway: gradient-aware sampling is
the better tool for last-layer predictive uncertainty, and conditioning---not the sampler alone---decides
whether the problem is tractable.
\end{abstract}

\section{Introduction}
Learned neural controllers produce a single trained parameter vector and treat its action predictions as
exact. On familiar states the controller is accurate; on states outside its training distribution it is
silently wrong, with no internal signal that anything has changed. For a controller driving a 7-DOF robot
arm through pick-and-place trajectories, this is a safety concern: the robot will confidently execute a
bad action rather than flag that it is operating outside its training regime.

The Bayesian view replaces the single parameter vector with a posterior distribution and predicts by
averaging over it. Familiar states produce tight action distributions; unfamiliar states produce
spread-out ones. That spread is the missing signal---an internal sense of ``I am not sure here'' that a
point-estimate controller cannot produce. Exact Bayesian inference in neural networks is intractable, so
in practice one relies on sampling-based approximations whose quality depends entirely on the MCMC method
behind them. This paper studies that dependence on a concrete controller, DNFC, taken from a prior
research project: a neural feedback controller that encodes a target (cartesian coordinates plus a
one-hot phase indicator) into a latent of the same shape as the robot state, subtracts the current state,
and passes the difference through a controller MLP that outputs $\tanh$-bounded joint velocities. We place
a Gaussian prior on the controller's final linear layer and compare two samplers head-to-head on the
resulting posterior: classical gradient-free RWMH and gradient-aware SGLD.

\textbf{Research question.} \emph{For Bayesian inference over the last layer of the DNFC controller, how do
RWMH and SGLD compare in terms of (i) mixing and effective sample size per unit of compute, and (ii) the
quality of the predictive uncertainty they produce on in-distribution, out-of-distribution
(extrapolation-split), and target-perturbed states?}

\textbf{Contributions.}
\begin{itemize}[leftmargin=1.5em,itemsep=1pt]
  \item A frozen-trunk Bayesian last-layer wrapper for a trained feedback controller, verified to
  reproduce the deterministic model exactly and to admit a correct tempered, subsampled log-posterior and
  gradient (Sec.~\ref{sec:foundation}).
  \item A rigorous RWMH-vs-SGLD comparison with split-$\hat R$/ESS diagnostics, three efficiency metrics
  that separate algorithm from hardware, and a closed-form Gaussian reference posterior, all averaged over
  five independently trained controllers (Secs.~\ref{sec:rwmh}--\ref{sec:uncertainty}).
  \item The empirical findings that (a) gradient information is necessary in this regime---RWMH is
  dimension-limited while SGLD is dimension-independent and $\sim\!33\times$ more compute-efficient;
  (b) \emph{conditioning} ($\sigma_a$), not the sampler, governs feasibility, and tempering does not help;
  and (c) last-layer MCMC yields a usable distribution-shift signal that SGLD recovers more faithfully
  than RWMH, while the dataset's nominal extrapolation split is not a genuine shift.
\end{itemize}

\section{Method}\label{sec:method}

\textbf{Model and Bayesian last layer.} DNFC encodes a target representation $g$ (9 cartesian coordinates
plus a 4-dimensional one-hot phase) through an MLP encoder into a latent of the same shape as the robot
state $s$, subtracts the state, and feeds the difference through a controller MLP that outputs
$\tanh$-bounded joint velocities (7 DOF). We freeze the encoder and all but the final linear layer of the
controller. Writing the frozen trunk's output as a feature $\phi(s,g)\in\mathbb{R}^{d}$, the action is
\[
f_\theta(s,g) = \tanh\!\big(W\,\phi(s,g) + b\big), \qquad \theta=(W,b)\in\mathbb{R}^{7(d+1)} .
\]
For the controller used here $d=384$, so the Bayesian parameter vector is
$\theta\in\mathbb{R}^{2695}$. Because the trunk is frozen, $\phi(s,g)$ is constant and is precomputed once
per state, making the posterior a cheap function of $\theta$ alone.

\textbf{Prior, likelihood, posterior.} A Gaussian prior $p(\theta)=\mathcal{N}(0,\alpha^2 I)$ is placed on
the last-layer weights, with $\alpha$ set to the empirical standard deviation of the trained weights
($\alpha\approx0.063$). The likelihood is Gaussian over actions with fixed observation noise $\sigma_a$,
\[
p(D\mid\theta)=\prod_{i=1}^{N'}\mathcal{N}\!\big(a_i \,\big|\, f_\theta(s_i,g_i),\;\sigma_a^2 I_7\big),
\qquad
p(\theta\mid D)\;\propto\;p(D\mid\theta)\,p(\theta),
\]
where $(s_i,g_i,a_i)$ is a (state, target, expert-action) triple from the DNFC dataset. Two controls keep
mixing feasible and are held identical across samplers: a tempered posterior $p(\theta\mid D)^{1/T}$ and
subsampling to a fixed subset of $N'$ training triples. We report $T$ and $N'$ up front
(Table~\ref{tab:setup}). \emph{Note on tempering.} Under $p(\theta\mid D)^{1/T}$ it is $T>1$ that flattens
the posterior; we use $T=1$ (untempered), because we find empirically (Sec.~\ref{sec:discussion}) that
tempering rescales posterior \emph{width} uniformly and therefore leaves the \emph{conditioning} that
governs mixing unchanged.

The single most consequential quantity is the posterior \emph{condition number}. Under a Gaussian
(Laplace) approximation it is
\[
\kappa \;=\; 1 + \frac{\alpha^2}{\sigma_a^2}\,\lambda_{\max}(\Phi^\top\Phi),
\]
where $\Phi$ is the matrix of precomputed features over the $N'$ subset. The feature Gram $\Phi^\top\Phi$
has \emph{effective rank} $\approx 3.6$ (the controller's input ``diff'' is only $14$-dimensional), so the
likelihood pins down a tiny subspace tightly and leaves the remaining $\sim\!2690$ directions to the
prior. This near-degeneracy is the source of the ill-conditioning that dominates our results.

\textbf{Samplers.} Both target the same posterior under the same $T$ and $N'$.
\emph{RWMH (gradient-free baseline)} uses isotropic Gaussian proposals
$\theta'=\theta+\tau\xi$, $\xi\sim\mathcal{N}(0,I)$, with $\tau$ tuned during burn-in (Robbins--Monro on
$\log\tau$) to the high-dimensional benchmark acceptance rate $\approx 23.4\%$.
\emph{SGLD (gradient-aware, primary)} uses minibatches of $B$ triples,
\[
\theta_{k+1}=\theta_k+\tfrac{\tau_k}{2}\Big(\tfrac{N'}{B}\!\!\sum_{i\in\mathcal{B}_k}\!\nabla_\theta\log
p(a_i\mid s_i,g_i,\theta)+\nabla_\theta\log p(\theta)\Big)+\sqrt{\tau_k}\,\xi_k,
\]
with the decreasing schedule $\tau_k\propto k^{-0.55}$ (floored to give a stationary sampling tail). We
implement the minibatch drift in closed form (no autograd graph) for speed.

\textbf{Analytic reference posterior.} Since the actions are small ($\sim\!10^{-2}$), $\tanh$ is
effectively the identity over their range, so each output dimension is a Bayesian linear regression on
the shared features and the posterior is a product of Gaussians with a single shared $385\times385$
covariance. This closed-form posterior is sampled i.i.d.\ and serves as an \emph{exact reference}: it
separates ``does the posterior encode a signal'' from ``did the sampler capture it.''

\textbf{Diagnostics.} For each sampler we run four chains from dispersed initial points and report
split-$\hat R$, ESS (Geyer initial-monotone estimator), and autocorrelation, with acceptance/step-size
trajectories. Because this small last layer is kernel-overhead-bound on a GPU (so minibatching buys no
wall-clock there), we report efficiency three ways to separate algorithm from hardware: ESS per
iteration, ESS per $10^6$ likelihood-term evaluations (the platform-independent ``per unit compute''
metric: RWMH touches $N'$ terms per step, SGLD touches $B$), and ESS per second of wall-clock.

\textbf{Predictive uncertainty.} From posterior samples $\{\theta^{(m)}\}_{m=1}^M$ we form the
Monte-Carlo predictive mean and spread,
\[
\bar\mu(s,g)=\tfrac1M\sum_m f_{\theta^{(m)}}(s,g),\qquad
\bar\sigma(s,g)^2=\tfrac1M\sum_m\big\|f_{\theta^{(m)}}(s,g)-\bar\mu(s,g)\big\|^2 .
\]
We evaluate three input regimes from the existing dataset: \emph{in-distribution} (interpolation split),
\emph{extrapolation} (extrapolation split), and \emph{target-perturbed} (in-distribution states with
Gaussian noise added to the target cartesian coordinates at several levels). We report action MSE against
the expert action, predictive-std distributions, extrapolation-detection AUROC (predictive std as the
score), and std-vs-error calibration, for the point estimate, the analytic reference, and the RWMH and
SGLD samples.

\textbf{Experimental setup.} Table~\ref{tab:setup} lists the configuration. To guard against
seed-specific flukes, all efficiency and uncertainty numbers are reported as mean~$\pm$~std across
\emph{five independently trained controllers} (different initialization seeds; identical architecture).
Runs use a single GPU; ESS estimates derive from four chains of thinned post-burn-in draws.

\begin{table}[h]
\centering\small
\begin{tabular}{@{}lll@{}}
\toprule
quantity & value & note \\
\midrule
Bayesian dimension $\theta$ & $2695$ & $d=384$ \\
prior scale $\alpha$ & $\approx 0.063$ & $=$ std of trained weights ($0.060$--$0.068$ across seeds) \\
observation noise $\sigma_a$ & $0.05$ & primary; train-residual $0.00127$ is too stiff (see below) \\
condition number $\kappa$ & $\approx 4\times10^{3}$ & vs $6.5\times10^{6}$ at the train-residual $\sigma_a$ \\
feature eff.\ rank & $\approx 3.6$ & $\lambda_{\max}(\Phi^\top\Phi)\approx 2.7\times10^{3}$ \\
subset / minibatch & $N'=4096$, $B=128$ & subset drawn from the train split \\
temperature & $T=1$ & untempered (tempering does not aid mixing) \\
chains / seeds & $4$ / $5$ & dispersed inits; non-AMC model, seeds $0$--$4$ \\
\bottomrule
\end{tabular}
\caption{Configuration. $\sigma_a$ sets the posterior conditioning $\kappa$ and is the key modelling knob.}
\label{tab:setup}
\end{table}

\section{Results}

\subsection{Foundation: the wrapper is exact and the posterior is correct}\label{sec:foundation}
The decomposed (frozen-trunk + Bayesian-head) forward pass reproduces the original monolithic controller
\emph{exactly} (maximum action difference $0$ in float32), and the trained point estimate explains the
training actions with $R^2=0.99$ on the $N'$ subset. The analytic log-posterior gradient matches central
finite differences to $\sim\!10^{-7}$, and the SGLD minibatch gradient (with the $N'/B$ rescaling) is an
unbiased estimator of the full-data gradient. These checks confirm that subsequent mixing behaviour
reflects the samplers, not implementation error.

\subsection{RWMH: correctly tuned, but defeated by dimension}\label{sec:rwmh}
RWMH adapts cleanly to the $23.4\%$ benchmark acceptance, yet the chains do not mix: on the (stiffest,
train-residual-$\sigma_a$) posterior the log-posterior split-$\hat R$ is $2.11$ with $\mathrm{ESS}=5.2$
from $8000$ draws, the per-coordinate $\hat R$ has median $6.2$ (max $20.4$), and \emph{no} coordinate
reaches $\hat R<1.1$. This is not an implementation defect: on isotropic Gaussians the sampler recovers
the target's mean and variance, and its ESS falls as $1/d$ (Table~\ref{tab:dimscale})---the classical
random-walk curse of dimension. The cause on the real posterior is anisotropy: the rank-$\approx\!3.6$
feature Gram yields a Gaussian-approximation condition number $\kappa\approx6.5\times10^{6}$, so isotropic
proposals limited by the steep directions cannot traverse the broad ones. Figure~\ref{fig:rwmh} shows the
traces, the adapted step size, and the flat autocorrelation.

\begin{table}[h]
\centering\small
\begin{tabular}{@{}lccc@{}}
\toprule
isotropic dimension $d$ & $100$ & $1000$ & $2695$ \\
\midrule
RWMH ESS (coordinate $0$) & $240$ & $13$ & $11$ \\
SGLD ESS (coordinate $0$) & $4221$ & $4775$ & $4554$ \\
\bottomrule
\end{tabular}
\caption{Mixing vs.\ dimension on isotropic Gaussian targets (same number of draws). RWMH degrades as
$1/d$; SGLD is essentially dimension-independent. This contrast is the mechanism behind every result below.}
\label{tab:dimscale}
\end{table}

\begin{figure}[t]
\centering
\includegraphics[width=0.96\textwidth]{rwmh/T1_N40000_C4/diagnostics.png}
\caption{RWMH diagnostics on the last-layer posterior: log-posterior and parameter traces, the
step-size adapted to $23.4\%$ acceptance, running acceptance, and the autocorrelation. The chains are
well tuned but do not mix.}
\label{fig:rwmh}
\end{figure}

\subsection{SGLD: gradient information pays off where it matters}\label{sec:sgld}
On isotropic targets SGLD mixes \emph{dimension-independently}---$\mathrm{ESS}\approx4554$ at $d=2695$
versus RWMH's $11$ (Table~\ref{tab:dimscale}), a $\sim\!400\times$ gap---because the gradient avoids the
random-walk slowdown. On the actual ill-conditioned posterior both samplers are limited by $\kappa$, so
\emph{per-iteration} mixing is comparable; SGLD's advantage is \emph{compute}, since it touches $B=128$
versus RWMH's $N'=4096$ likelihood terms per step (Table~\ref{tab:efficiency},
Fig.~\ref{fig:efficiency}). Averaged over five seeds, SGLD is $32.9\times\!\pm\!15.4$ more efficient per
unit of compute, and the per-seed ratio exceeds one on every seed ($9.3$--$56\times$). The wall-clock
caveat is honest: on this small last layer the GPU is kernel-overhead-bound, so RWMH's full-batch step is
as cheap as SGLD's minibatch ($0.14$ vs $0.15$\,ms) and RWMH is marginally ahead in ESS/second; on a CPU
(FLOP-bound) SGLD's step is $2.3\times$ cheaper, so the minibatch advantage is real and simply does not
surface in GPU wall-clock at this scale. A single-model sweep over $\sigma_a$ confirms SGLD's
per-compute lead is stable across conditioning ($27$--$95\times$).

\begin{table}[h]
\centering\small
\begin{tabular}{@{}lccc@{}}
\toprule
efficiency metric (log-posterior) & RWMH & SGLD & SGLD/RWMH \\
\midrule
ESS / iteration & $(3.2\pm1.5)\times10^{-4}$ & $(2.7\pm0.8)\times10^{-4}$ & $\approx 1$ \\
\textbf{ESS / $10^{6}$ term-evals} & $0.078\pm0.037$ & $2.11\pm0.63$ & $\mathbf{32.9\times\!\pm\!15.4}$ \\
ESS / second (GPU) & $2.26\pm1.1$ & $1.75\pm0.52$ & $0.8$ (overhead-bound) \\
\bottomrule
\end{tabular}
\caption{Sampler efficiency at $\sigma_a=0.05$, mean~$\pm$~std over 5 seeds. Per-iteration mixing is
comparable (both $\kappa$-limited); SGLD's edge is per unit compute.}
\label{tab:efficiency}
\end{table}

\begin{figure}[t]
\centering
\includegraphics[width=0.92\textwidth]{multiseed/multiseed_efficiency.png}
\caption{RWMH vs.\ SGLD efficiency across five trained seeds. Left: log-posterior ESS per $10^6$
term-evaluations (per-seed points; bars are mean~$\pm$~std). Right: the per-seed SGLD/RWMH compute-
efficiency ratio, above parity for every seed.}
\label{fig:efficiency}
\end{figure}

\subsection{Predictive uncertainty}\label{sec:uncertainty}
\textbf{Genuine input shift is detected, and SGLD recovers it best.} Perturbing the target coordinates
raises the predictive std monotonically, and the sampler ordering is analytic $\gtrsim$ SGLD $>$ RWMH at
every level and on every seed (Table~\ref{tab:uncertainty}, Fig.~\ref{fig:uncertainty}). RWMH
systematically under-estimates the spread---its poor mixing shrinks the apparent variance---while SGLD
tracks the exact analytic posterior closely. This is the project's positive result: last-layer MCMC
produces a usable distribution-shift signal, and the gradient-aware sampler reproduces it far more
faithfully.

\textbf{The nominal extrapolation split is not out-of-distribution.} Action error does not rise on the
``extrapolation'' split (MSE ratio extrap/interp $=0.84\pm0.04$), and predictive std fails to separate it
from in-distribution states---detection AUROC $\approx0.5$ \emph{even for the exact analytic posterior}
($0.460\pm0.010$; Table~\ref{tab:uncertainty}). This is a correct null, not a method failure: $0.0\%$ of
the extrapolation split's target coordinates fall outside the training range (interpolation itself is
$18.2\%$ outside), and feature-space $z$-scores against training are identical ($0.60$ vs $0.61$). The
frozen encoder maps these inputs on-manifold, the controller generalizes, and the uncertainty correctly
reports ``not more uncertain.'' We therefore use the controlled target-perturbation experiment as the
distribution-shift demonstration and report the extrapolation-split null transparently.

\begin{table}[h]
\centering\small
\begin{tabular}{@{}lccc@{}}
\toprule
method & median in-dist.\ std & extrap.\ AUROC & perturb.\ std @ $0.2$\,m \\
\midrule
analytic (exact reference) & $0.0186$ & $0.460\pm0.010$ & $0.0927\pm0.0089$ \\
SGLD & $0.0167$ & $0.457\pm0.046$ & $0.0834\pm0.0104$ \\
RWMH & $0.0147$ & $0.475\pm0.034$ & $0.0584\pm0.0056$ \\
point estimate & $0$ & $0.5$ (no spread) & $0$ \\
\bottomrule
\end{tabular}
\caption{Predictive uncertainty, mean~$\pm$~std over 5 seeds. AUROC $\approx0.5$ everywhere (the extrap
split is not OOD), but predictive std rises with genuine perturbation, with SGLD closest to the exact
reference. Point-estimate action MSE extrap/interp $=0.84\pm0.04$ ($\approx1$: extrap no harder).}
\label{tab:uncertainty}
\end{table}

\begin{figure}[t]
\centering
\includegraphics[width=0.92\textwidth]{multiseed/multiseed_uncertainty.png}
\caption{Predictive uncertainty across five seeds. Left: predictive std vs.\ target-perturbation level
(mean $\pm$ std bands); the rise is monotonic and SGLD tracks the exact analytic posterior while RWMH
lags. Right: extrapolation-detection AUROC per seed, at chance for all methods including the exact
reference---the nominal extrapolation split is not a genuine shift.}
\label{fig:uncertainty}
\end{figure}

\section{Discussion}\label{sec:discussion}
\textbf{Gradient information is necessary in this regime.} The $1/d$ scaling of RWMH versus the
dimension-independent mixing of SGLD (Table~\ref{tab:dimscale}) is the root cause of every downstream
difference. At $\theta\in\mathbb{R}^{2695}$ a gradient-free random walk simply cannot explore the
posterior, whereas SGLD can---and does so while touching far less data per step.

\textbf{Conditioning, not the sampler, governs feasibility.} The observation noise $\sigma_a$ sets
$\kappa=1+(\alpha^2/\sigma_a^2)\lambda_{\max}(\Phi^\top\Phi)$. At the ``data-driven'' train-residual
$\sigma_a=0.00127$ the posterior is a needle ($\kappa\approx6.5\times10^{6}$) that neither sampler can
explore \emph{and} whose predictive spread is near zero---useless as an uncertainty signal. A deliberate
$\sigma_a=0.05$ ($\kappa\approx4\times10^{3}$) restores both samplability and a meaningful signal.
Crucially, \emph{tempering does not substitute for this}: $p(\theta\mid D)^{1/T}$ rescales every posterior
direction by the same factor, leaving $\kappa$---and hence mixing---unchanged (we verified $T{=}100$
mixes no better than $T{=}1$). The lever that matters is the likelihood precision relative to the prior,
i.e.\ $\sigma_a$.

\textbf{SGLD is the better tool for last-layer predictive uncertainty.} Beyond raw efficiency, SGLD's
samples reproduce the exact posterior's spread and its response to perturbation, whereas RWMH's poor
mixing makes it under-report uncertainty---a failure mode that is invisible without a reference but
directly harmful for a safety signal. The closed-form analytic posterior was essential precisely because
it disentangles the posterior's content from the sampler's fidelity.

\textbf{On the research question.} The outcome is the ``interesting trade-off'' the question anticipated:
SGLD dominates on dimension scaling and per-unit-compute efficiency and on uncertainty fidelity, while the
two are comparable per iteration and (on this small problem) in GPU wall-clock. And the uncertainty study
yields a usable extrapolation signal only under genuine shift; the dataset's labelled extrapolation split
is not one---a caveat that is itself a useful, honest finding about evaluating learned controllers.

\section{Limitations and Future Work}
$\sigma_a=0.05$ is a reported modelling choice rather than a learned quantity; its effect on conditioning
is mapped by a single-model sweep but not propagated through the uncertainty experiments. The MCMC ESS is
small in absolute terms (tens of effective samples), so the RWMH/SGLD predictive stds carry estimator
noise; the analytic reference is the reliable ground truth here. Finally, the absence of a genuine
out-of-hull extrapolation split prevented a true OOD-detection AUROC; constructing such a split (target
positions outside the training convex hull), propagating $\sigma_a$ sensitivity through the predictive
metrics, and running a closed-loop uncertainty rollout in simulation are natural next steps.

\section{Conclusion}
We made the last layer of a trained neural feedback controller Bayesian and showed how the choice of MCMC
sampler shapes the resulting uncertainty. A well-tuned gradient-free random walk is correct but
dimension-limited and mixes poorly at $\theta\in\mathbb{R}^{2695}$; gradient-aware SGLD mixes
dimension-independently, is $\sim\!33\times$ more efficient per unit of compute (robust across five
trained seeds), and recovers the exact posterior's predictive uncertainty far more faithfully. Whether
the problem is tractable at all is decided by posterior conditioning---set by the observation-noise scale,
not by tempering. The Bayesian last layer does provide a usable ``I am not sure here'' signal under
genuine input shift; the dataset's nominal extrapolation split simply is not such a shift. For equipping a
point-estimate controller with calibrated uncertainty, gradient-aware sampling is the better tool, and the
modelling of the likelihood precision is as important as the sampler itself.

\section*{Reproducibility}
All experiments are reproducible with the \texttt{bayes/} package (see \texttt{bayes/README.md} and
\texttt{bayes/RESULTS.md}): \texttt{verify\_week1} (foundation), \texttt{run\_rwmh}, \texttt{run\_sgld},
\texttt{run\_week4} (predictive uncertainty), and \texttt{run\_multiseed} (the five-seed aggregation that
produces Tables~\ref{tab:efficiency}--\ref{tab:uncertainty} and Figs.~\ref{fig:efficiency}--\ref{fig:uncertainty}).

\end{document}
